\chapter[USB TMC Driver]{USB TMC Driver}
\label{cp:USBTMC-Driver}

\section*{Resumen}
Estoy usando la librería de ESP-IDF para implementar el driver USB TMC. Esta librería permite la comunicación con dispositivos USB a través de la interfaz TMC (Test and Measurement Class) pero no esta completamente implementada en el framework de ESP-IDF. La implementación se basa en TinyUSB, una librería de código abierto para la comunicación USB en microcontroladores en la que si esta completamente implementada la Clase TMC. La implementación tambien incluye la creacion de una interfaz estandar para darle al motor SCPI un aapi estandarizada para la comunicación con el driver USB TMC y cualquier otro que se quiera implementar en el futuro.

\section{Implementación del driver USB TMC en ESP-IDF}
Como en la librería de ESP-IDF no esta completamente implementada la Clase TMC, debemos agregar a mano los descriptores de la clase TMC y agregarlos con la función de inicialización
\begin{minted}{c}
    /**
    * @brief Install the TinyUSB device driver and start the TinyUSB task.
    *
    * This helper configures the USB PHY when requested, prepares descriptors,
    * initializes the TinyUSB stack, and starts the TinyUSB task.
    *
    * @note When supplying a custom composite device descriptor with an Interface
    *       Association Descriptor, keep `bDeviceClass` as `TUSB_CLASS_MISC` and
    *       `bDeviceSubClass` as `MISC_SUBCLASS_COMMON`.
    *
    * @param[in] config TinyUSB stack configuration.
    *
    * @return
    *      - ESP_OK on success
    *      - ESP_ERR_INVALID_ARG if `config` is NULL or contains unsupported port or task settings
    *      - ESP_ERR_INVALID_STATE if the TinyUSB device task is already running
    *      - ESP_ERR_NO_MEM if memory allocation fails during startup
    *      - Other error codes from TinyUSB task startup, USB PHY setup, or descriptor setup
    */
    esp_err_t tinyusb_driver_install(const tinyusb_config_t *config);
\end{minted}

\section{Descriptores de la clase TMC}


\subsection{Descriptor del dispositivo}
\begin{minted}{c}
    tusb_desc_device_t const desc_device = {
        .bLength = sizeof(tusb_desc_device_t),
        .bDescriptorType = TUSB_DESC_DEVICE,
        .bcdUSB = 0x0200,
        .bDeviceClass = TUD_USBTMC_CLASS,           // 0x01 para descriptor de de tipo dispositivo
        .bDeviceSubClass = TUD_USBTMC_SUBCLASS,     // 0xFE
        .bDeviceProtocol = TUD_USBTMC_PROTO_488,    // 0x03
        .bMaxPacketSize0 = CFG_TUD_ENDPOINT0_SIZE,  // 64 bytes para maximo tamaño de paquete de control

        .idVendor = USB_VID,                        // 0x16C0
        .idProduct = USB_PID,                       // 0x2044
        .bcdDevice = USB_BCD,                       // 0x0100

        .iManufacturer = 0x01, // Índice del string del fabricante
        .iProduct = 0x02,      // Índice del string del producto
        .iSerialNumber = 0x03, // Índice del string del número de serie
        .bNumConfigurations = 0x01};
\end{minted}

\subsection{Descriptor de configuración}

\subsection{Descriptor de interfaz}
\subsection{Descriptopres de cadenas de caracteres}

\section{Callbacks weaks implementadios en TinyUSB para acceso al stack de USB}
Tenemos que al menos implementar los callbacks de TinyUSB para poder acceder al stack de USB y poder enviar y recibir datos a traves de la interfaz TMC. Estos callbacks son:
\begin{minted}{c}
    // 2. Callback de apertura (Requerido para inicializar el bus)*
    void tud_usbtmc_open_cb(uint8_t interface_id);
    // 3. Callback de recepción de mensajes BULK
    bool tud_usbtmc_msgBulkOut_start_cb(usbtmc_msg_request_dev_dep_out const *msgHeader)
    // 4. Callback de datos recibidos
    bool tud_usbtmc_msg_data_cb(void *data, size_t len, bool transfer_complete);
    // 5. Callback de petición BULK IN (El host pide datos)*
    bool tud_usbtmc_msgBulkIn_request_cb(usbtmc_msg_request_dev_dep_in const *request);
    // 6. Callback de finalización BULK IN*
    bool tud_usbtmc_msgBulkIn_complete_cb(void);
    // 7. Callback para el STATUS BYTE (IEEE 488) - CRITICO PARA SCPI
    uint8_t tud_usbtmc_get_stb_cb(uint8_t *tmcResult);
    // 8. Callback para el TRIGGER (Ejemplo: comando *TRG o señal GET)
    bool tud_usbtmc_msg_trigger_cb(usbtmc_msg_generic_t *msg);
    // 9. Clear Feature Bulk OUT
    void tud_usbtmc_bulkOut_clearFeature_cb(void);
    // 10. Clear Feature Bulk IN
    void tud_usbtmc_bulkIn_clearFeature_cb(void);
    // 11. Initiate Abort Bulk OUT
    bool tud_usbtmc_initiate_abort_bulk_out_cb(uint8_t *tmcResult);
    // 12. Check Abort Bulk OUT
    bool tud_usbtmc_check_abort_bulk_out_cb(usbtmc_check_abort_bulk_rsp_t *rsp);
    // 13. Initiate Abort Bulk IN
    bool tud_usbtmc_initiate_abort_bulk_in_cb(uint8_t *tmcResult);
    // 14. Check Abort Bulk IN
    bool tud_usbtmc_check_abort_bulk_in_cb(usbtmc_check_abort_bulk_rsp_t *rsp);
    // 15. Initiate Clear
    bool tud_usbtmc_initiate_clear_cb(uint8_t *tmcResult);
    // 16. Check Clear
    bool tud_usbtmc_check_clear_cb(usbtmc_get_clear_status_rsp_t *rsp);
\end{minted}

\section{Máquina de estados para la comunicación}
La maquina de estados es minimalista, tiene 4 estados principales: 
\begin{minted}{c}
    typedef enum
    {
        STATE_TMC_IDLE,
        STATE_TMC_RECEIVING,
        STATE_TMC_PROCESSING,
        STATE_TMC_REPLY_READY
    } usb_tmc_state_t;
\end{minted}

Esta maquina puede controlar 5 eventos del host y 1 del SCPI:
\begin{minted}{c}
    typedef enum
    {
        EV_TMC_RX_START,   // Inicia recepción (Bulk OUT)
        EV_TMC_RX_CHUNK,   // Llega un pedazo de datos
        EV_TMC_RX_END,     // Llega el final del mensaje
        EV_TMC_TX_REQ,     // Host pide leer (Bulk IN)
        EV_TMC_TX_DONE,    // Host terminó de leer
        EV_TMC_SCPI_DONE   // libscpi no encontró query
    } usb_tmc_event_t;
\end{minted}

\subsection{STATE\_TMC\_IDLE}
\subsubsection*{Eventos controlados permitidos}
En el primer estado \textbf{STATE\_TMC\_IDLE}, hay solo 2 eventos que pueden ocurrir: \textbf{EV\_TMC\_RX\_START} y \textbf{EV\_TMC\_TX\_REQ}. 
Cuando el host envía un mensaje, se genera el evento \textbf{EV\_TMC\_RX\_START} y la FSM pasa al estado \textbf{STATE\_TMC\_RECEIVING}. 
Cuando el host solicita leer, se genera el evento \textbf{EV\_TMC\_TX\_REQ}, en general el host solo puede pedir datos al equipo si antes hizo alguna consulta y esto pone la maquina de estados en \textbf{STATE\_TMC\_PROCESSING}.

\subsubsection*{Eventos controlados no permitidos}
Hacer bulkIN (pedido de datos) sin antes hacer una consulta genera un evento \textbf{Error -420: Query Unterminated"} que debemos informar a libscpi para que lo registre. Luego responderemos con un NAK/Empty al host para que continue sin esperar la respuesta.

\subsection{STATE\_TMC\_RECEIVING}
\subsubsection*{Eventos controlados permitidos}
En este estado solo existen 3 eventos posibles: \textbf{EV\_TMC\_RX\_CHUNK} y \textbf{EV\_TMC\_RX\_END} y \textbf{EV\_TMC\_TX\_REQ}.
Despues de iniciada una recepción los datos pueden llegar en una sola transferencia, si entran en el endpoint del usb que tiene un tamaño de 64 bytes, o pueden llegar en varios pedazos si el mensaje es mas largo. 
Cada vez que llega un pedazo de datos se genera el evento \textbf{EV\_TMC\_RX\_CHUNK} y la FSM permanece en el estado \textbf{STATE\_TMC\_RECEIVING}. 
Cuando llega el final del mensaje se genera el evento \textbf{EV\_TMC\_RX\_END} y la FSM pasa al estado \textbf{STATE\_TMC\_PROCESSING}.

\subsubsection*{Eventos controlados no permitidos}
En este caso la unica forma de generar un evento no permitido es que el host pida datos antes de que termine de enviar el mensaje, pero esto solo se produce por un error del controlador usb en el host dispositivo, ya que el host no puede enviar un pedido de datos mientras esta enviando un mensaje.

\subsection{STATE\_TMC\_PROCESSING}
\subsubsection{Eventos controlados permitidos}
Aqui solo hay un evento posible: \textbf{EV\_TMC\_SCPI\_DONE}, que se genera cuando libscpi termina de procesar el mensaje recibido. En este mensaje pueden venir desde ninguna a muchas consultas así que la maquina de estados deve revisar el stream de salida para verificar si hay datos disponibles y la FSM pasa al estado \textbf{STATE\_TMC\_REPLY\_READY} o si no hay datos disponivles volver a \textbf{STATE\_TMC\_IDLE}.

\subsubsection{Eventos controlados no permitidos}
En este estado el motor scpi puede tardar algo de tiempo en procesar el mensaje recibido. Si el host envia una nueva consulta generardo un evento de tipo \textbf{EV\_TMC\_RX\_START} antes de que libscpi termine de procesar el mensaje, se genera un evento \textbf{Error -410: Query Interrupted} que debemos informar a libscpi para que lo registre. Tambien se debe reconfigurar la FSM volviendo al estado \textbf{STATE\_TMC\_RECEIVING} para recibir el nuevo mensaje. Tambien se debe descartar cualquier respuesta que esté generando libscpi, lo que se hace limpiando el buffer del libscpi con una bandera atomica antes de enviar el evento \textbf{EV\_TMC\_SCPI\_DONE}.

\subsection{STATE\_TMC\_REPLY\_READY}
\subsubsection{Eventos controlados permitidos}
En este estado hay 2 eventos permitidos posibles: \textbf{EV\_TMC\_TX\_REQ} y \textbf{EV\_TMC\_TX\_DONE}. Cuando el host solicita leer, se genera el evento \textbf{EV\_TMC\_TX\_REQ} y la FSM envia los datos al host. Cuando el host termina de leer, se genera el evento \textbf{EV\_TMC\_TX\_DONE} y la FSM vuelve al estado \textbf{STATE\_TMC\_IDLE}.

\subsubsection{Eventos controlados no permitidos}
Aqui nuevamente se puede generar un evento no permitido si el host envia una nueva consulta generardo un evento de tipo \textbf{EV\_TMC\_RX\_START} antes de que el host termine de leer la respuesta. Esto genera un evento \textbf{Error -410: Query Interrupted} que debemos informar a libscpi para que lo registre. Tambien se debe reconfigurar la FSM volviendo al estado \textbf{STATE\_TMC\_RECEIVING} para recibir el nuevo mensaje. tambien es importante abortar la transmision de datos al host. Este evento no deberia ser posible si el host esta bien implementado, ya que no deberia enviar una nueva consulta mientras esta leyendo la respuesta a la consulta anterior.

